\subsubsection{谱兼容外置、半直积修复与有限/无穷二影刚性}\label{subsubsec:xi-spectrum-compatible-semidirect-two-shadow-rigidity}
本段在 Lee--Yang 四层覆盖的既有单值化框架上，将“可交换外置障碍—非交换修复—线性椭圆双影刚性”组织为同一条可检验链条。

\begin{theorem}[可交换外置寄存器对非交换单值群的不可完备性]\label{thm:xi-terminal-abelian-register-obstruction-nonabelian-monodromy}
设 \(U\) 为道路连通拓扑空间，\(\pi:X\to U\) 为连通有限分支覆盖，分支集记为 \(B\subset U\)，并取基点 \(u_0\in U\setminus B\)。
令
\[
\rho:\pi_1(U\setminus B,u_0)\longrightarrow \mathfrak S_n
\]
为单值表示，且假设 \(\rho(\pi_1(U\setminus B,u_0))\) 非交换。

称三元组 \((A,\iota,D)\) 为该覆盖的一组可交换外置寄存器数据，若：
\begin{enumerate}
  \item \(A\) 为交换群；
  \item \(\iota:\pi_1(U\setminus B,u_0)\to A\) 为群同态；
  \item \(D:\{1,\dots,n\}\times A\to\{1,\dots,n\}\) 满足对任意 \(\gamma\in\pi_1(U\setminus B,u_0)\) 与 \(i\in\{1,\dots,n\}\),
  \[
  \rho(\gamma)(i)=D\bigl(i,\iota(\gamma)\bigr).
  \]
\end{enumerate}
则此类三元组不存在。
\end{theorem}
\begin{proof}
由 \(\rho(\pi_1(U\setminus B,u_0))\) 非交换，可取 \(\alpha,\beta\in\pi_1(U\setminus B,u_0)\) 使
\[
\rho([\alpha,\beta])\neq \mathrm{id},
\qquad
[\alpha,\beta]:=\alpha\beta\alpha^{-1}\beta^{-1}.
\]
由 \(A\) 交换且 \(\iota\) 为同态，
\[
\iota([\alpha,\beta])
=\iota(\alpha)\iota(\beta)\iota(\alpha)^{-1}\iota(\beta)^{-1}
=e_A
=\iota(1).
\]
故对任意 \(i\),
\[
D\bigl(i,\iota([\alpha,\beta])\bigr)=D(i,e_A)=D\bigl(i,\iota(1)\bigr).
\]
左侧按定义等于 \(\rho([\alpha,\beta])(i)\)，右侧等于 \(\rho(1)(i)=i\)，从而 \(\rho([\alpha,\beta])(i)=i\) 对所有 \(i\) 成立，矛盾于 \(\rho([\alpha,\beta])\neq\mathrm{id}\)。
证毕。
\end{proof}

\begin{corollary}[Lee--Yang 四层覆盖中的可交换外置障碍]\label{cor:xi-terminal-abelian-register-obstruction-leyang-s4}
在定理 \ref{thm:xi-terminal-zm-leyang-elliptic-structure} 给出的 Lee--Yang 四层覆盖 \(\pi_y:E_{\min}\to\mathbb P^1_y\) 上，
\[
\mathrm{Mon}(\pi_y)\cong S_4
\]
为非交换群，故满足定理 \ref{thm:xi-terminal-abelian-register-obstruction-nonabelian-monodromy} 的不可完备性结论：任何要求路径拼接同态化的可交换外置寄存器都不能实现全局单值恢复。
\end{corollary}

\begin{theorem}[位移半直积 Gödel 寄存器对非交换算法轨迹的忠实外置]\label{thm:xi-terminal-semidir-godel-register-faithful}
设 \(\Sigma\) 为有限字母表（\(|\Sigma|\ge 2\)），\(\Sigma^\ast\) 为自由幺半群（串联为乘法），并取单射编码
\(
c:\Sigma\hookrightarrow \NN_{\ge 1}
\)。
记 \((p_i)_{i\ge 1}\) 为递增素数序列。

对 \(w=a_1\cdots a_n\in\Sigma^\ast\) 定义位序 Gödel 编码
\[
G(w):=\prod_{i=1}^{n}p_i^{\,c(a_i)},
\qquad
|w|=n.
\]
定义素数位移同态
\[
\mathrm{Sh}^{k}\!\left(\prod_{i\ge1}p_i^{e_i}\right):=\prod_{i\ge1}p_{i+k}^{e_i}.
\]
令
\[
\mathcal R:=\NN\times \NN_{>0},
\qquad
(n,A)\odot(m,B):=\bigl(n+m,\ A\cdot \mathrm{Sh}^{n}(B)\bigr).
\]
则映射
\[
J:\Sigma^\ast\to\mathcal R,\qquad
J(w):=(|w|,G(w))
\]
为单射幺半群同态，且 \((|w|,G(w))\) 可唯一恢复 \(w\)。
此外，对任意交换幺半群 \(M\)，不存在单射幺半群同态 \(\Sigma^\ast\hookrightarrow M\)。
\end{theorem}
\begin{proof}
\(\mathrm{Sh}^{n}\circ\mathrm{Sh}^{m}=\mathrm{Sh}^{n+m}\) 且每个 \(\mathrm{Sh}^n\) 为幺半群同态，故 \(\odot\) 结合，单位元为 \((0,1)\)。
对 \(u=a_1\cdots a_n\)、\(v=b_1\cdots b_m\)，
\[
G(uv)
=\left(\prod_{i=1}^{n}p_i^{c(a_i)}\right)\left(\prod_{j=1}^{m}p_{n+j}^{c(b_j)}\right)
=G(u)\,\mathrm{Sh}^{n}(G(v)),
\]
故 \(J(uv)=J(u)\odot J(v)\)。

单射性由唯一分解定理给出：\(J(w)=(n,G(w))\) 已给出长度 \(n\)，且
\(
v_{p_i}(G(w))=c(a_i)
\)
可逐位恢复 \(a_i\)（因 \(c\) 单射）。

若 \(f:\Sigma^\ast\to M\) 为到交换幺半群的同态，取 \(a\neq b\in\Sigma\)，则
\[
f(ab)=f(a)f(b)=f(b)f(a)=f(ba).
\]
若 \(f\) 单射则推出 \(ab=ba\)，与自由幺半群非交换性矛盾。证毕。
\end{proof}

\input{sections/generated/eq_xi_prime_register_semidirect_two_shadow_audit}

\begin{theorem}[行列式素因子分解给出的椭圆面积守恒律]\label{thm:xi-terminal-ellipse-area-euler-prime-law}
设 \(A\in\mathrm{GL}_2(\QQ)\)，单位圆盘 \(D:=\{x\in\RR^2:\|x\|\le 1\}\)，椭圆盘 \(E_A:=A(D)\)。
则
\[
\mathrm{Area}(E_A)=\pi\,|\det A|.
\]
并且若
\(
\det A=\pm\prod_p p^{v_p(\det A)}
\)
（仅有限个 \(v_p\neq 0\)），则
\[
\mathrm{Area}(E_A)=\pi\prod_p p^{v_p(\det A)}.
\]
\end{theorem}
\begin{proof}
二维 Lebesgue 测度在可逆线性变换下按雅可比因子 \( |\det A| \) 缩放，故
\(
\mathrm{Area}(A(D))=|\det A|\,\mathrm{Area}(D)=\pi|\det A|
\)。
其余结论由有理数乘法群的素因子分解直接得到。证毕。
\end{proof}

\begin{theorem}[有理椭圆合同类的二元不变量完备性]\label{thm:xi-terminal-rational-ellipse-two-invariant-completeness}
设 \(A\in\mathrm{GL}_2(\QQ)\)，记 \(E_A:=A(S^1)\subset\RR^2\)。
定义
\[
\Delta_1(A):=(\det A)^2\in\QQ_{>0},
\qquad
\Delta_2(A):=\mathrm{tr}(A^\top A)\in\QQ_{>0}.
\]
则 \(E_A\) 的欧氏合同类由 \((\Delta_1(A),\Delta_2(A))\) 完全决定。
若 \(\sigma_1\ge \sigma_2>0\) 为 \(A\) 的奇异值，则
\[
\sigma_1^2,\sigma_2^2\ \text{是}\ t^2-\Delta_2(A)t+\Delta_1(A)=0\ \text{的两根},
\]
因而
\[
\sigma_{1,2}
=\sqrt{\frac{\Delta_2(A)\pm\sqrt{\Delta_2(A)^2-4\Delta_1(A)}}{2}}.
\]
特别地，对任意 \(A,B\in\mathrm{GL}_2(\QQ)\),
\[
E_A\cong E_B\ (\text{欧氏合同})
\Longleftrightarrow
\bigl(\Delta_1(A),\Delta_2(A)\bigr)=\bigl(\Delta_1(B),\Delta_2(B)\bigr).
\]
\end{theorem}
\begin{proof}
极分解 \(A=UP\)（\(U\) 正交，\(P=(A^\top A)^{1/2}\) 对称正定）给出
\(
E_A=U(P(S^1))
\)；
故合同类由 \(P\) 决定。
\(P\) 的特征值为 \(\sigma_1,\sigma_2\)，即椭圆半轴长。
又
\[
\sigma_1^2+\sigma_2^2=\mathrm{tr}(A^\top A)=\Delta_2(A),
\qquad
\sigma_1^2\sigma_2^2=\det(A^\top A)=(\det A)^2=\Delta_1(A),
\]
因此 \((\sigma_1,\sigma_2)\) 由 \((\Delta_1,\Delta_2)\) 唯一恢复，反之亦然。证毕。
\end{proof}

\input{sections/generated/tab_xi_elliptic_two_shadow_invariants_audit}

\begin{theorem}[有限素部格与实部椭圆影的二影刚性]\label{thm:xi-terminal-finite-archimedean-two-shadow-rigidity}
定义 prime-register 纤维
\[
\widehat{\ZZ}:=\varprojlim_m \ZZ/m\ZZ\cong\prod_p \ZZ_p.
\]
对 \(A\in\mathrm{GL}_2(\QQ)\)，记有限素部格影
\[
\Lambda_A:=A\,\widehat{\ZZ}^{\,2}\subset \mathbb A_f^{\,2},
\]
并记实部椭圆影 \(E_A:=A(S^1)\subset\RR^2\)。
若 \(A,B\in\mathrm{GL}_2(\QQ)\) 满足
\[
\Lambda_A=\Lambda_B,
\qquad
E_A=E_B,
\]
则存在 \(Q\in O_2(\ZZ)\) 使 \(A=BQ\)。
若再施加定向保持规范（等价于 \(\det(B^{-1}A)=1\)），则 \(Q\in SO_2(\ZZ)\)，因此仅余 \(4\) 个离散候选；
进一步施加任一破缺该四元对称的规范锚定后，\(Q\) 唯一，从而 \(A\) 由 \((\Lambda_A,E_A)\) 唯一确定。
\end{theorem}
\begin{proof}
令 \(C:=B^{-1}A\)。
由 \(\Lambda_A=\Lambda_B\) 得
\(
C\,\widehat{\ZZ}^{\,2}=\widehat{\ZZ}^{\,2}
\)，故 \(C\in\mathrm{GL}_2(\widehat{\ZZ})\)。

由 \(E_A=E_B\) 得同一椭圆二次型：
\[
x^\top(AA^\top)^{-1}x=1
\Longleftrightarrow
x^\top(BB^\top)^{-1}x=1,
\]
从而 \(AA^\top=BB^\top\)。
左乘 \(B^{-1}\)、右乘 \((B^{-1})^\top\) 得
\[
CC^\top=I,
\]
即 \(C\in O_2(\QQ)\)。

同时 \(C\in\mathrm{GL}_2(\widehat{\ZZ})\) 蕴含其每个有理条目在所有 \(\ZZ_p\) 中整，故条目属于
\(
\QQ\cap\bigcap_p\ZZ_p=\ZZ
\)；
即 \(C\in\mathrm{GL}_2(\ZZ)\)。
综上 \(C\in O_2(\ZZ)\)，即 \(A=BQ\)（\(Q:=C\in O_2(\ZZ)\)）。

当 \(\det Q=1\) 时 \(Q\in SO_2(\ZZ)\)。
二维整正交群中 \(|SO_2(\ZZ)|=4\)，故仅有四个定向保持候选；对任意固定规范锚定（例如长轴正向锚定）即可消去残余离散歧义并得到唯一性。证毕。
\end{proof}

